\chapter{Continuum States, Quantized Momentum, and Interference}

\section{When a Discrete Label Becomes Continuous}

Begin with an observable whose possible values are labeled by integers
\(n=1,2,\ldots\). Its probabilities obey
\begin{equation}
  \sum_n P(n)=1,
  \label{eq:discrete-normalization}
\end{equation}
and the corresponding basis vectors may be chosen orthonormally:
\begin{equation}
  \braket m n=\delta_{mn}.
  \label{eq:kronecker}
\end{equation}
The Kronecker delta is one when its two labels agree and zero otherwise. It
can equally well be regarded as a function of the difference \(m-n\).

\begin{figure}[tbp]
  \centering
  \includegraphics[width=0.54\textwidth]{lecture_04_figure_02.png}
  \caption{Discrete probability normalization and basis orthonormality.}
  \lectureframe{00:03:11}
\end{figure}

For a state \(\ket\psi\), the probability of obtaining \(n\) is
\begin{equation}
  P(n)=\left|\braket n\psi\right|^{2}.
  \label{eq:discrete-born}
\end{equation}
These three statements are the pattern we must preserve when the label is no
longer discrete.

Let the measured quantity now be the position \(x\) of a particle on a line.
For an ordinary continuous distribution, a single point has zero
probability. The useful object is therefore a probability density
\(\rho(x)\). The probability of finding the particle in a short interval is
\begin{align}
  \operatorname{Prob}\!\left(x_1\leq x\leq x_2\right)
  &=\int_{x_1}^{x_2}\rho(x)\,dx,\\
  \int_{-\infty}^{\infty}\rho(x)\,dx&=1.
  \label{eq:continuous-probability}
\end{align}
Thus sums become integrals, and probability masses become densities.

\begin{figure}[tbp]
  \centering
  \includegraphics[width=0.58\textwidth]{lecture_04_figure_03.png}
  \caption{For a continuous coordinate, probability belongs to an interval
  and is represented by area under a density.}
  \lectureframe{00:05:59}
\end{figure}

\section{The Dirac Delta}

The continuum analogue of the Kronecker delta is the Dirac delta
distribution. Its informal picture is an infinitely narrow spike at
\(x=y\), zero away from that point, with unit area:
\begin{equation}
  \delta(x-y)=0\quad (x\neq y),
  \qquad
  \int_{-\infty}^{\infty}\delta(x-y)\,dx=1.
  \label{eq:delta-unit-area}
\end{equation}
No ordinary function has these properties. The precise definition is its
action on a sufficiently regular test function:
\begin{equation}
  \int_{-\infty}^{\infty}\delta(x-y)f(x)\,dx=f(y).
  \label{eq:l04-delta-sampling}
\end{equation}
The delta is therefore best understood as a distribution under an integral,
not as a function with a literal infinite value.

\begin{figure}[tbp]
  \centering
  \includegraphics[width=0.62\textwidth]{lecture_04_figure_04.png}
  \caption{A narrow unit-area spike samples the value of a smooth function at
  \(x=y\).}
  \lectureframe{00:12:00}
\end{figure}

A useful approximation makes the spike picture explicit. Let
\(\delta_\epsilon(x-y)\) be a rectangle of width \(\epsilon\) and height
\(1/\epsilon\), centered at \(y\). Its area is one, and for continuous \(f\),
\begin{equation}
  \int \delta_\epsilon(x-y)f(x)\,dx
    \longrightarrow f(y)
  \qquad (\epsilon\to0).
\end{equation}
The unit area must not be confused with a unit Hilbert-space norm:
\(\int\left|\delta_\epsilon\right|^{2}dx=1/\epsilon\) diverges in the limit.

Position basis states therefore obey delta normalization:
\begin{equation}
  \braket x y=\delta(x-y),
  \qquad
  \int_{-\infty}^{\infty}dx\,\ket x\bra x=\mathbf 1.
  \label{eq:position-basis}
\end{equation}
Distinct positions are orthogonal. The formal quantity
\(\braket x x=\delta(0)\) is divergent because \(\ket x\) is a generalized
eigenstate, not a normalizable vector in the Hilbert space.

\section{The Position Amplitude and Born Density}

The abstract probability rule has not changed. In the discrete case,
\(\left|\braket n\psi\right|^{2}\) is the probability for the outcome \(n\). In the
continuous case,
\begin{equation}
  \rho(y)=\left|\braket y\psi\right|^{2}
  \label{eq:continuum-born}
\end{equation}
is the probability density at \(y\). To see what this means in the position
representation, write
\begin{equation}
  \psi(x):=\braket x\psi.
\end{equation}
The wavefunction of the generalized position state \(\ket y\) is
\(\braket x y=\delta(x-y)\). Hence
\begin{align}
  \braket y\psi
    &=\int dx\,\braket yx\braket x\psi \\
    &=\int dx\,\delta(x-y)\psi(x)
      =\psi(y).
  \label{eq:position-component}
\end{align}
Born's rule is therefore
\begin{align}
  \rho(y)&=\left|\psi(y)\right|^{2},\\
  \operatorname{Prob}(y\in[a,b])
  &=\int_a^b\left|\psi(y)\right|^{2}\,dy.
  \label{eq:born-density}
\end{align}

\begin{figure}[tbp]
  \centering
  \includegraphics[width=0.62\textwidth]{lecture_04_figure_05.png}
  \caption{Delta sampling identifies the position amplitude
  \(\braket y\psi\) with the wavefunction value \(\psi(y)\).}
  \lectureframe{00:22:00}
\end{figure}

\begin{classroomqa}
\audiencequestion{Could the first probability postulate be stated once more
for the continuous case?\lecturetimestamp{00:22:39--00:26:37}}
\lecturerresponse{For a discrete observable, the probability of outcome
\(n\) is \(\left|\braket n\psi\right|^{2}\). For position, the same expression
\(\left|\braket y\psi\right|^{2}\) is a density. Since
\(\braket y\psi=\psi(y)\), the probability of an interval is the integral of
\(\left|\psi(y)\right|^{2}\) over that interval. Dirac's notation was designed so
that the discrete and continuous statements have the same form.}
\end{classroomqa}

\section{A Lattice Explains Continuum Normalization}

The divergent self-overlap of \(\ket x\) becomes transparent if the line is
temporarily replaced by a lattice with spacing \(a\). Let
\(x_n=na\), and let \(\ket n\) be ordinary normalized site states:
\begin{equation}
  \braket n m=\delta_{nm},
  \qquad
  \sum_n\ket n\bra n=\mathbf 1.
\end{equation}
Define continuum-normalized states at the lattice points by
\begin{equation}
  \ket{x_n}:=\frac{1}{\sqrt a}\ket n.
  \label{eq:lattice-rescaling}
\end{equation}
Their overlap and completeness relations are
\begin{equation}
  \braket{x_n}{x_m}=\frac{\delta_{nm}}{a},
  \qquad
  \sum_n a\,\ket{x_n}\bra{x_n}=\mathbf 1.
  \label{eq:lattice-delta}
\end{equation}
As \(a\to0\), the sum \(\sum_n a\) becomes an integral, while
\(\delta_{nm}/a\) becomes \(\delta(x-y)\). The height \(1/a\) and width \(a\)
again give unit area.

\begin{figure}[tbp]
  \centering
  \includegraphics[width=0.61\textwidth]{lecture_04_figure_06.png}
  \caption{Rescaling lattice kets by \(a^{-1/2}\) produces Dirac-delta
  normalization in the continuum limit.}
  \lectureframe{00:30:00}
\end{figure}

\begin{classroomqa}
\audiencequestion{How can the overlap of a state with itself be infinite
without ruining the probability interpretation?
\lecturetimestamp{00:26:38--00:34:24}}
\lecturerresponse{The position eigenstates are not physical normalized
states; they are generalized basis vectors. On a lattice their normalized
ancestors are perfectly ordinary. The factor \(a^{-1/2}\) converts a discrete
probability into a density, and it converts \(\delta_{nm}\) into
\(\delta_{nm}/a\). Finite wave packets remain normalized, and every delta
appears under an integral where its unit-area property gives a finite
answer.}
\end{classroomqa}

\section{A Particle on a Circle}

Now constrain the particle to a circle of radius \(R\). Cutting the circle at
one point turns it into the interval
\begin{equation}
  0\leq x<2\pi R,
\end{equation}
but the two ends represent the same physical point. The wavefunction must
therefore satisfy
\begin{equation}
  \psi(x+2\pi R)=\psi(x).
  \label{eq:periodic-boundary}
\end{equation}
Periodic functions form a complex vector space: scalar multiples and sums
of periodic functions are still periodic.

\begin{figure}[tbp]
  \centering
  \includegraphics[width=0.60\textwidth]{lecture_04_figure_07.png}
  \caption{Cutting the circle gives an interval whose wavefunction values
  must match after one circumference.}
  \lectureframe{00:38:00}
\end{figure}

On this space the momentum operator is
\begin{equation}
  \hat p=-i\hbar\frac{d}{dx},
  \label{eq:l04-momentum-operator}
\end{equation}
with periodic boundary conditions included in its domain. Its eigenvalue
equation,
\begin{equation}
  -i\hbar\frac{d\psi_p}{dx}=p\psi_p,
\end{equation}
has the plane-wave solution
\begin{equation}
  \psi_p(x)=C e^{ipx/\hbar}.
  \label{eq:plane-wave-circle}
\end{equation}
The name momentum will acquire its full physical justification only after
time evolution and the motion of wave packets are developed. For now it is
the eigenvalue of the translation generator.

\begin{figure}[tbp]
  \centering
  \includegraphics[width=0.54\textwidth]{lecture_04_figure_08.png}
  \caption{The momentum eigenfunction is a complex plane wave.}
  \lectureframe{00:41:40}
\end{figure}

Normalization on the finite circle fixes \(C\):
\begin{align}
  1
    &=\int_0^{2\pi R}\left|\psi_p(x)\right|^{2}\,dx
      =2\pi R\left|C\right|^{2},\\
  \psi_p(x)
    &=\frac{1}{\sqrt{2\pi R}}e^{ipx/\hbar}.
  \label{eq:normalized-plane-wave}
\end{align}
The constant magnitude means that a momentum eigenstate has uniform position
density around the circle.

\section{Periodicity Quantizes Momentum}

The normalized plane wave is not automatically an allowed state. It must
also obey the periodic boundary condition:
\begin{align}
  e^{ip(x+2\pi R)/\hbar}&=e^{ipx/\hbar},\\
  e^{i2\pi Rp/\hbar}&=1.
\end{align}
The phase gained in one circuit must be an integer multiple of \(2\pi\):
\begin{equation}
  \frac{2\pi Rp}{\hbar}=2\pi n,
  \qquad n\in\mathbb Z.
\end{equation}
Thus
\begin{equation}
  p_n=\frac{n\hbar}{R},
  \qquad
  \psi_n(x)=\frac{1}{\sqrt{2\pi R}}e^{inx/R}.
  \label{eq:circle-spectrum}
\end{equation}

\begin{figure}[tbp]
  \centering
  \includegraphics[width=0.58\textwidth]{lecture_04_figure_09.png}
  \caption{Single-valuedness around the circle restricts momentum to
  \(p_n=n\hbar/R\).}
  \lectureframe{00:50:00}
\end{figure}

The allowed functions are mutually orthogonal:
\begin{equation}
  \int_0^{2\pi R}\psi_m^*(x)\psi_n(x)\,dx=\delta_{mn}.
\end{equation}
Here position is continuous while momentum is discrete. In this problem the
discreteness is a consequence of compactness and periodicity; discrete
spectra in other systems can arise from other boundary conditions or from a
confining potential.

The spacing between neighboring momenta is
\begin{equation}
  \Delta p=\frac{\hbar}{R}.
\end{equation}
As \(R\to\infty\), this spacing tends to zero. A continuum-normalized momentum
basis may be obtained by the same rescaling used for position:
\begin{align}
  \ket{p_n}_{\mathrm c}
  &:=\frac{1}{\sqrt{\Delta p}}\ket n,\\
  {}_{\mathrm c}\!\braket{p_m}{p_n}_{\mathrm c}
  &=\frac{\delta_{mn}}{\Delta p}\\
  &\longrightarrow\delta(p-p').
\end{align}

For orbital motion on the circle, the angular momentum is
\begin{equation}
  L=Rp,
\end{equation}
and therefore
\begin{equation}
  L_n=Rp_n=n\hbar.
  \label{eq:angular-momentum}
\end{equation}
The radius cancels because angular momentum measures phase winding around an
angle. This conclusion applies to a scalar orbital wavefunction with
\(\psi(\phi+2\pi)=\psi(\phi)\). Half-integer spin belongs to a different
representation and is a later subject.

\begin{classroomqa}
\audiencequestion{Does \(R\to\infty\) recover classical mechanics?
\lecturetimestamp{00:56:58--00:57:20}}
\lecturerresponse{No. It only makes the momentum spectrum dense. Classical
behavior requires more than closely spaced eigenvalues; in particular, it
requires suitable wave packets and a dynamical limit in which position and
momentum can both be narrow compared with the relevant macroscopic scales.}
\end{classroomqa}

\begin{classroomqa}
\audiencequestion{Why does \(R\) remain in the momentum spacing but disappear
from the angular-momentum spectrum?
\lecturetimestamp{00:57:37--00:58:48}}
\lecturerresponse{A fixed phase winding \(n\) around a larger circle has a
smaller tangential wave number and hence smaller \(p_n=n\hbar/R\).
Multiplying by the lever arm \(R\) restores \(L_n=n\hbar\). The quantized
object is the action associated with a full angular circuit.}
\end{classroomqa}

\section{Compatibility and Common Eigenstates}

A classical particle may be assigned a position and a momentum at the same
instant. Quantum mechanically, a position eigenstate is localized like a
delta distribution, whereas a momentum eigenstate is a plane wave with
uniform position density. These states cannot be the same.

The algebraic question is whether two Hermitian operators \(\hat A\) and
\(\hat B\) possess a complete common eigenbasis. If
\begin{equation}
  \hat A\ket n=a_n\ket n,
  \qquad
  \hat B\ket n=b_n\ket n,
\end{equation}
then
\begin{equation}
  \hat A\hat B\ket n=a_nb_n\ket n
  =\hat B\hat A\ket n.
\end{equation}
Define the commutator
\begin{equation}
  [\hat A,\hat B]:=\hat A\hat B-\hat B\hat A.
\end{equation}
Compatible observables commute. Conversely, commuting Hermitian matrices can
be simultaneously diagonalized, including within degenerate eigenspaces.
For unbounded operators, the same statement requires compatible domains and
the corresponding spectral conditions; those analytic details are
suppressed here.

\section{The Canonical Commutator}

For position and momentum,
\begin{equation}
  (\hat x\psi)(x)=x\psi(x),
  \qquad
  (\hat p\psi)(x)=-i\hbar\psi'(x).
\end{equation}
Apply their commutator to an arbitrary differentiable wavefunction:
\begin{align}
  [\hat x,\hat p]\psi
    &=x\left(-i\hbar\frac{d\psi}{dx}\right)
      +i\hbar\frac{d}{dx}\bigl(x\psi\bigr)\\
    &=-i\hbar x\psi'
      +i\hbar\bigl(\psi+x\psi'\bigr)\\
    &=i\hbar\psi.
\end{align}
Since this holds for every state in the common domain,
\begin{equation}
  [\hat x,\hat p]=i\hbar\,\mathbf 1.
  \label{eq:canonical-commutator}
\end{equation}

\begin{figure}[tbp]
  \centering
  \includegraphics[width=0.66\textwidth]{lecture_04_figure_10.png}
  \caption{The product rule supplies the uncancelled term in
  \([\hat x,\hat p]\psi=i\hbar\psi\).}
  \lectureframe{01:13:00}
\end{figure}

The commutator is nonzero on every nonzero state. If a vector were a common
eigenvector of \(\hat x\) and \(\hat p\), their commutator would annihilate
it, contradicting \eqref{eq:canonical-commutator}. Position and momentum
therefore have no common eigenvector. The uncertainty relation will later
quantify this incompatibility.

\begin{classroomqa}
\audiencequestion{Is every state an eigenvector of the commutator, and can
\(i\hbar\) be regarded as an observable value?
\lecturetimestamp{01:15:33--01:16:45}}
\lecturerresponse{Every vector is indeed an eigenvector of
\([\hat x,\hat p]=i\hbar\mathbf 1\), with the same eigenvalue. But the
commutator of two Hermitian operators is anti-Hermitian, so it is not itself
an observable. The Hermitian combination
\(-i[\hat x,\hat p]=\hbar\mathbf 1\) contains no state-dependent measurement
information. Its importance is that it fixes the algebra relating position
and momentum.}
\end{classroomqa}

\begin{classroomqa}
\audiencequestion{Why is \(x\psi(x)\) not an eigenvalue times
\(\psi(x)\)?\lecturetimestamp{01:16:51--01:19:07}}
\lecturerresponse{An eigenvalue is one scalar constant for the whole vector.
The coordinate \(x\) varies across the function. A generic Gaussian, for
example, is therefore not an eigenfunction of the position operator.
Only a generalized state supported at one point satisfies
\(x\delta(x-x_0)=x_0\delta(x-x_0)\).}
\end{classroomqa}

\section{One Aperture and Two Alternatives}

The two-slit experiment turns the linear structure of states into a visible
probability pattern. An incoming particle meets a barrier with two small
openings, and a screen records the position \(y\) at which it arrives.
With only one opening available, the outgoing wave spreads. In three
dimensions a schematic far-field contribution has the form
\begin{equation}
  \psi(r)\sim\frac{e^{ikr}}{r},
  \qquad
  r=\sqrt{\ell^2+y^2},
  \label{eq:outgoing-wave}
\end{equation}
where \(\ell\) is the distance from aperture to screen. Over a small region
of the screen it may be approximated by a local plane wave.

\begin{figure}[tbp]
  \centering
  \includegraphics[width=0.64\textwidth]{lecture_04_figure_11.png}
  \caption{A small opening produces an outgoing wave whose local phase
  gradient changes across the screen.}
  \lectureframe{01:23:00}
\end{figure}

Now open both holes. Let \(\psi_1(y)\) and \(\psi_2(y)\) be the amplitudes
associated with the two alternatives. The state-space rule is
\begin{equation}
  \psi(y)=\psi_1(y)+\psi_2(y),
  \label{eq:add-amplitudes}
\end{equation}
not the addition of the separate probabilities. Locally, suppressing slowly
varying envelopes and temporarily setting \(\hbar=1\), write
\begin{equation}
  \psi_1=e^{ipy},
  \qquad
  \psi_2=e^{iqy}.
\end{equation}
Then
\begin{align}
  \left|\psi_1+\psi_2\right|^{2}
    &=(e^{ipy}+e^{iqy})(e^{-ipy}+e^{-iqy})\\
    &=2+e^{i(p-q)y}+e^{-i(p-q)y}\\
    &=2+2\cos\bigl((p-q)y\bigr).
  \label{eq:local-interference}
\end{align}
The cross terms retain the relative phase. They create alternating bright
and dark fringes even though either local plane wave separately has constant
magnitude.

\begin{figure}[tbp]
  \centering
  \includegraphics[width=0.62\textwidth]{lecture_04_figure_12.png}
  \caption{The cross terms in the squared sum generate a nonnegative
  oscillating detection density.}
  \lectureframe{01:30:00}
\end{figure}

\begin{figure}[tbp]
  \centering
  \begin{tikzpicture}[scale=0.9,>=Latex]
    \draw[thick] (0,-2.2) -- (0,-0.45);
    \draw[thick] (0,0.45) -- (0,1.15);
    \draw[thick] (0,1.65) -- (0,2.2);
    \fill (-3,0.6) circle (1.5pt) node[left] {source};
    \draw[->] (-2.9,0.6) -- (-0.18,0.1);
    \draw[->] (-2.9,0.6) -- (-0.18,1.4);
    \draw[thick] (5,-2.3) -- (5,2.3);
    \draw (0,0.1) -- (5,1.55);
    \draw (0,1.4) -- (5,1.55);
    \fill (5,1.55) circle (1.5pt) node[right] {\(y\)};
    \draw[<->] (-0.35,0.1) -- (-0.35,1.4)
      node[midway,left] {\(d\)};
    \node[below] at (0,-2.2) {barrier};
    \node[below] at (5,-2.3) {screen};
    \node at (2.55,-1.55) {\(\ell\)};
  \end{tikzpicture}
  \caption{The relative path length from two separated openings changes with
  the screen coordinate.}
\end{figure}

In the far field, if the slit separation is \(d\), the path difference at
angle \(\theta\) is approximately \(d\sin\theta\). For wavelength
\(\lambda=2\pi/k\), adjacent bright fringes satisfy
\begin{equation}
  d\sin\theta_m=m\lambda.
\end{equation}
At small angles, \(y\approx\ell\theta\), so
\begin{equation}
  \Delta y\approx\frac{\lambda\ell}{d}.
  \label{eq:fringe-spacing}
\end{equation}
Moving the holes farther apart makes the fringes closer; moving them together
spreads the fringes apart.

\begin{classroomqa}
\audiencequestion{What happens in the local model if \(p=q\)?
\lecturetimestamp{01:32:25--01:33:08}}
\lecturerresponse{Then the relative phase is constant:
\(\psi_1+\psi_2=2e^{ipy}\). The intensity is the constant \(4\), in the
chosen normalization. There are no alternating fringes, but the two
contributions are everywhere fully constructive.}
\end{classroomqa}

\begin{classroomqa}
\audiencequestion{Is this the Fraunhofer diffraction pattern?
\lecturetimestamp{01:38:39--01:39:57}}
\lecturerresponse{It is the far-field two-aperture part of that pattern. The
local beat model isolates the rapidly varying two-path fringes. A full
Fraunhofer calculation also retains the finite width of each opening, which
multiplies those fringes by the broader one-aperture diffraction envelope.}
\end{classroomqa}

\section{Which-Path Information Must Wait}

Suppose a detector at the barrier records which opening the electron used.
The detector must then become correlated with the alternative. Schematically,
\begin{equation}
  \ket\Psi
    =\ket{\text{path }1}\ket{D_1}
     +\ket{\text{path }2}\ket{D_2}.
  \label{eq:path-detector}
\end{equation}
When the detector records are distinguishable, the electron alone no longer
has the same interference term. A full explanation requires composite
systems and entanglement, so it belongs after the tensor-product structure
has been developed.

\begin{classroomqa}
\audiencequestion{Why does a detector placed at the openings destroy the
interference pattern?\lecturetimestamp{01:40:04--01:41:53}}
\lecturerresponse{Detecting the path means leaving different imprints in
another physical system. The electron and detector must then be described as
one composite state, with each path correlated with a different detector
record. The interference disappears when those records are distinguishable.
The mechanism is entanglement, whose derivation must wait until composite
systems are available.}
\end{classroomqa}

\section{Phase Gradients and Elastic Deflection}

For a wave of the form \(e^{i\phi(y)}\), the local vertical momentum is the
phase gradient
\begin{equation}
  p_y(y)=\hbar\frac{d\phi}{dy}.
\end{equation}
Near the central axis the distance \(r=\sqrt{\ell^2+y^2}\) changes slowly
with \(y\), so the vertical phase gradient is small. Farther from the axis,
the wave has a larger vertical momentum component. This does not require an
increase in kinetic energy. In an elastic scattering picture, the magnitude
of the momentum remains fixed:
\begin{equation}
  p_x^2+p_y^2=p^2.
\end{equation}
An increase in \(\left|p_y\right|\) is accompanied by a decrease in
\(\left|p_x\right|\); the momentum vector changes direction rather than length.

\begin{classroomqa}
\audiencequestion{Where does the energy come from when the particle acquires
vertical momentum?\lecturetimestamp{01:41:59--01:44:12}}
\lecturerresponse{It need not acquire additional energy. The barrier can
redirect the momentum elastically, much as a stationary wall changes the
direction of a bouncing ball. The outgoing wavelength, and therefore the
magnitude of the momentum, can remain the same while its horizontal and
vertical components change.}
\end{classroomqa}

\section{Particles Arrive One at a Time}

The wavefunction describes both alternatives, but each detection event is a
localized spot. Send one electron at a time and the screen records one spot
at a time. Repeat the preparation many times and the accumulated spots
approach the density \(\left|\psi_1+\psi_2\right|^{2}\). The density cannot be
negative because it is an absolute square, but it can vanish where the two
amplitudes cancel.

\begin{classroomqa}
\audiencequestion{Is it then meaningless to speak of a corpuscular particle
in this experiment?\lecturetimestamp{01:44:19--01:45:35}}
\lecturerresponse{No. The particle aspect is visible in each localized
detection. The wave aspect is visible in the probability law governing many
such detections. A single event is a spot; the interference pattern is the
statistical arrangement of many spots.}
\end{classroomqa}

\begin{classroomqa}
\audiencequestion{How does this compare with the interference of
light?\lecturetimestamp{01:45:39--01:47:17}}
\lecturerresponse{The same distinction appears with photons. An intense beam
draws the fringe pattern quickly. If the beam is attenuated enough, the
screen registers separated photon events, yet those events accumulate into
the same interference distribution. Reducing the flux changes the number of
particles present at once, not the underlying probability law.}
\end{classroomqa}

\section{The One-Aperture Envelope}

A single opening does not give a featureless amplitude. A convenient local
form is
\begin{equation}
  \psi_1(y)=\rho(y)e^{i\phi(y)},
  \label{eq:one-hole-envelope}
\end{equation}
where \(\rho(y)\) changes slowly while the phase may change rapidly. Its
detection density is
\begin{equation}
  \left|\psi_1(y)\right|^2=\rho(y)^2.
\end{equation}
The rapid phase cancels against its complex conjugate, leaving a smooth
envelope. With a second displaced opening,
\begin{equation}
  \psi_2(y)=\rho_2(y)e^{i\phi_2(y)},
\end{equation}
and the probability becomes
\begin{equation}
  \left|\psi_1+\psi_2\right|^{2}
  =\rho_1^2+\rho_2^2
   +2\rho_1\rho_2\cos(\phi_1-\phi_2).
  \label{eq:enveloped-interference}
\end{equation}
The smooth one-aperture envelopes are modulated by the relative-phase
fringes.

\begin{figure}[tbp]
  \centering
  \includegraphics[width=0.62\textwidth]{lecture_04_figure_13.png}
  \caption{A rapidly varying phase multiplied by a smooth envelope gives a
  smooth one-aperture probability density.}
  \lectureframe{01:49:00}
\end{figure}

\begin{classroomqa}
\audiencequestion{If the amplitude oscillates behind one opening, why does
the screen show only a broad blob?\lecturetimestamp{01:47:28--01:51:36}}
\lecturerresponse{The oscillation is in complex phase. Multiplication by the
complex conjugate removes that phase, leaving the slowly varying
\(\rho(y)^2\). With two openings, a relative phase remains in the cross term,
and that is what produces fringes.}
\end{classroomqa}

A finite aperture can also diffract by itself. Different points across a
sharp opening contribute different phases, so a single-slit envelope may
contain maxima and minima of its own. This is distinct from the finer
two-slit modulation, although a real experiment displays both at once.

\begin{classroomqa}
\audiencequestion{Can the two edges of one sharp opening create their own
interference pattern?\lecturetimestamp{01:51:37--01:52:41}}
\lecturerresponse{Yes. A finite opening is a continuum of secondary
contributions, and its edges make the phase structure especially visible.
The resulting single-aperture diffraction determines the broad envelope;
the separation of two openings determines the additional two-path fringes.}
\end{classroomqa}

\section{Visibility, Wavelength, and Idealization}

Exact dark fringes require coherent contributions of equal magnitude. If the
two apertures have unequal amplitudes \(A_1\) and \(A_2\), then
\begin{equation}
  I(y)=\left|A_1\right|^{2}+\left|A_2\right|^{2}
       +2\left|A_1A_2\right|\cos\Delta\phi(y).
\end{equation}
The minimum is
\((\left|A_1\right|-\left|A_2\right|)^2\), so unequal illumination fills in the dark
fringes. Noise, finite source coherence, detector resolution, and imperfect
geometry do the same. Under ideal conditions the minima can nevertheless be
extremely dark.

\begin{classroomqa}
\audiencequestion{Are there screen positions that will never be hit, even
when particles are sent one at a time?
\lecturetimestamp{01:52:43--01:53:43}}
\lecturerresponse{In the ideal equal-amplitude model, yes: destructive
interference makes the probability exactly zero at the dark fringes. Real
apparatus introduces background, unequal amplitudes, and finite resolution,
so observed minima are generally small rather than mathematically perfect.}
\end{classroomqa}

\begin{classroomqa}
\audiencequestion{Do unequal hole sizes matter, and how small must the
openings be?\lecturetimestamp{01:53:44--01:55:54}}
\lecturerresponse{Small, similarly illuminated openings give the clearest
two-path pattern. Unequal sizes alter the two amplitudes and reduce fringe
visibility. The relevant scales are the opening widths, their separation,
the propagation distance, and the de Broglie wavelength
\(\lambda=h/p\). Long-wavelength particles diffract through larger
structures; short-wavelength electrons require correspondingly finer
geometry.}
\end{classroomqa}

Passing through a sufficiently small opening prepares an approximate
outgoing state centered at that opening. In this limited sense the later
pattern can become insensitive to many details of the source. Coherence
still matters: the source and apparatus must not carry distinguishable
which-path information.

\begin{classroomqa}
\audiencequestion{Does the particle forget its source after passing through
the small opening?\lecturetimestamp{01:55:55--01:56:41}}
\lecturerresponse{The opening can prepare a new approximate outgoing state,
so many earlier details cease to affect the downstream shape. That is an
idealization, not complete independence from the source: momentum spread,
coherence, and correlations with other systems can still matter.}
\end{classroomqa}

\section{Reuse Is Not Reversal}

The final question separates two ideas that are easy to confuse. One may
collect an electron, prepare it again, and repeat the experiment with that
same particle; repeated trials do not require newly created electrons.
Reversing an entire detected event is a different matter. The screen and all
other correlated degrees of freedom would have to be included in the
reversal.

\begin{classroomqa}
\audiencequestion{Could the same electron be sent through repeatedly, and is
that the same as reversing the experiment?
\lecturetimestamp{01:56:42--01:59:16}}
\lecturerresponse{The same electron may be recovered, re-prepared, and used
for another trial, producing the same statistics as any identically prepared
electron. Reversing a completed detection is not the same operation. To
decide whether the full evolution can be reversed, one must first introduce
Hamiltonians, time evolution, and the detector as part of the quantum
system.}
\end{classroomqa}

That is where the calculation properly stops. The kinematics has supplied
continuous bases, delta normalization, quantized momentum on a circle,
noncommuting position and momentum, and interference from added amplitudes.
The next step is dynamics: a Hamiltonian must say how these states evolve in
time.
